\documentclass[11pt,a4paper]{article}
\usepackage[utf8]{inputenc}\usepackage[T1]{fontenc}
\usepackage[provide=*,english]{babel}
\usepackage{lmodern,microtype,amsmath,amssymb,booktabs,geometry,xcolor,fancyhdr}
\usepackage[hidelinks]{hyperref}
\geometry{margin=2.3cm,headheight=14pt}
\setlength{\parindent}{0pt}\setlength{\parskip}{0.4em}
\definecolor{gcblue}{RGB}{34,73,110}\definecolor{gcgray}{RGB}{90,96,102}
\pagestyle{fancy}\fancyhf{}
\fancyhead[L]{\small\color{gcgray}GC2D: ABBA4 single projection}
\fancyhead[R]{\small\color{gcgray}Model theory}\fancyfoot[C]{\thepage}
\newcommand{\R}{\mathbb{R}}
\begin{document}
\begin{center}
 {\LARGE\bfseries\color{gcblue}ABBA4 Implicit Single Projection}\par
 \vspace{0.35em}{\large One symmetric projection around an unprojected triple jump}
\end{center}

\section{Unprojected fourth-order base}

Let $\mathcal A_{q,t}$ be one signed endpoint-time ABBA map on two copies,
with no intermediate diagonal projection.  Define
\[
\gamma=\frac{1}{2-\sqrt[3]{2}},\qquad
\delta=-\frac{\sqrt[3]{2}}{2-\sqrt[3]{2}}.
\]
The complete base map is
\[
\mathcal C_{h,t}=
\mathcal A_{\gamma h,t+(\gamma+\delta)h}
\circ\mathcal A_{\delta h,t+\gamma h}
\circ\mathcal A_{\gamma h,t}.
\]
Because $2\gamma+\delta=1$ and $2\gamma^3+\delta^3=0$, the palindromic
composition cancels the leading third-order defect of a symmetric second-order
base.  The middle duration is negative, and the two copies pass continuously
through all three factors.

\section{One outer symmetric projection}

For
\[
E=\begin{pmatrix}I\\I\end{pmatrix},\quad
P=\tfrac12\begin{pmatrix}I&I\end{pmatrix},\quad
G=\begin{pmatrix}I&-I\end{pmatrix},\quad N=G^T,
\]
the reduced formulation is
\begin{align}
M(\mu)&=\mathcal C_{h,t}(Ez+N\mu),\\
r(\mu)&=GM(\mu)+2\mu=0,\\
z^+&=P(M(\mu)+N\mu).
\end{align}
If $J_j$ is the Jacobian of signed factor $j$, then
$J_\mathcal C=J_3J_2J_1$ in execution order and
\[
D_\mu r=GJ_\mathcal C N+2I.
\]
Every residual evaluation traverses all three ABBA factors.

\section{Configurations}

The class exposes the same orthogonal axes as the other implicit ABBA methods:
two projection formulations, Newton or Broyden, and physical, shared-time, or
fully extended state.  The selector values are global to the outer step.  In
the simultaneous formulation the corrected doubled output joins $\mu$ in the
unknown; elimination gives the reduced residual above.  State extension
changes what is duplicated, but never introduces intermediate projection.

\section{Difference from \texttt{ABBA4Implicit}}

Both fourth-order classes use $(\gamma,\delta,\gamma)$, but
\texttt{ABBA4Implicit} solves three projections and re-embeds a physical state
between signed factors.  This method solves one projection around the full
unprojected composition.  Since projection and factor evolution do not
commute, the classes define different numerical maps rather than two costs for
the same map.

\section{Properties and limitations}

The ideal method is symmetric and designed for order four.  Geometric claims
refer to the converged projection root; finite nonlinear tolerances perturb the
map.  Negative internal time increments require the potential and dynamics to
be evaluable slightly outside the outer interval.  Observer records preserve
the three ordered factors so the accepted-map tangent can be reconstructed.

\section{Implementation correspondence}

The public class and outer residual are in
\path{src/simulation/methods/abba/order4_implicit_single_projection.py}.
The common ABBA stages and projection helpers live in the neighboring private
modules.  The companion architecture document under \texttt{../simulation/}
gives dimensions, diagnostics, and runtime flow.

\end{document}
